% Figure: contour plot of f(x,y) = x^2 + 2 y^2 with gradient arrows
% perpendicular to the contours, illustrating the steepest-ascent
% theorem. Used in section 11.2 prose.
\begin{figure}[htbp]
\centering
\begin{tikzpicture}[scale=1.0,
  contour/.style={engblue, thick},
  arrow/.style={->, engred, very thick},
  axis/.style={engblack, thick, ->},
  label/.style={font=\small},
]
% Axes
\draw[axis] (-3.2, 0) -- (3.4, 0) node[right, label] {$x$};
\draw[axis] (0, -2.4) -- (0, 2.6) node[above, label] {$y$};

% Concentric ellipses: f(x,y) = x^2 + 2 y^2 = c, so x^2/c + y^2/(c/2) = 1
% Semi-axes (sqrt c, sqrt(c/2)). Choose c = 0.5, 1, 2, 4.
\draw[contour] (0,0) ellipse ({sqrt(0.5)} and {sqrt(0.25)});
\draw[contour] (0,0) ellipse ({sqrt(1.0)} and {sqrt(0.5)});
\draw[contour] (0,0) ellipse ({sqrt(2.0)} and {sqrt(1.0)});
\draw[contour] (0,0) ellipse ({sqrt(4.0)} and {sqrt(2.0)});

% Contour labels
\node[label, engblue] at (2.15, 0.18) {$c=4$};
\node[label, engblue] at (1.50, 0.18) {$c=2$};
\node[label, engblue] at (1.05, 0.18) {$c=1$};

% Gradient arrows at a few sample points: grad f = (2x, 4y).
% Scale arrows for visual clarity (length = 0.15 * grad).
% Point A: (1, 0.5). grad = (2, 2). Length 0.25.
\draw[arrow] (1, 0.5) -- ++({2*0.25}, {2*0.25});
\filldraw[engred] (1, 0.5) circle (0.04);
% Point B: (-1, 0.7). grad = (-2, 2.8).
\draw[arrow] (-1, 0.7) -- ++({-2*0.25}, {2.8*0.25});
\filldraw[engred] (-1, 0.7) circle (0.04);
% Point C: (1.4, -0.6). grad = (2.8, -2.4).
\draw[arrow] (1.4, -0.6) -- ++({2.8*0.25}, {-2.4*0.25});
\filldraw[engred] (1.4, -0.6) circle (0.04);
% Point D: (0, 1.0). grad = (0, 4).
\draw[arrow] (0, 1.0) -- ++(0, {4*0.18});
\filldraw[engred] (0, 1.0) circle (0.04);

% Annotation: gradient is perpendicular to the contour at every point
\node[label, anchor=west] at (1.55, 1.1) {$\grad f \perp$ contour};

\end{tikzpicture}
\caption[Contours of $f(x,y) = x^{2} + 2y^{2}$ with gradient
arrows.]{Contours of $f(x, y) = x^{2} + 2y^{2}$ (blue ellipses) with
$\grad f = (2x, 4y)$ drawn at four sample points (red). At every
point, the gradient is perpendicular to the contour through that
point, and its magnitude reports the rate of steepest ascent. The
asymmetry of the ellipses reflects the asymmetric coefficient on
$y^{2}$: the function increases twice as fast per unit of $y$ as
per unit of $x$.}
\label{fig:vol02:ch11:contour-gradient}
\end{figure}
